\section{The Mass--Radius Relationship}
\label{sec:mass-radius}

The TOV integration of Section~\ref{sec:tov}, carried out across a range
of central energy densities $\varepsilon_c$ with the free neutron Fermi
gas equation of state, produces a continuous curve in the mass-radius
plane. This mass-radius ($M$--$R$) curve is the primary observable
prediction of the model and the natural point at which theory meets
observation.
% ------------------------------------------------------------
\subsection{Structure of the Mass--Radius Curve}
% ------------------------------------------------------------

The $M$--$R$ curve for the free neutron Fermi gas has a characteristic
shape that reflects the interplay between the equation of state and
gravity as the central density is varied. Beginning at low central
densities, both $M$ and $R$ are small. As $\varepsilon_c$ increases, the
mass rises steeply while the radius varies more slowly, broadly
consistent with the Newtonian scaling $R \propto M^{-1/3}$ derived in
Section~\ref{sec:newtonian}: the star becomes more massive and more
compact simultaneously. The curve reaches a maximum mass at
$M_{\mathrm{OV}} \approx 0.7\,M_\odot$, corresponding to a radius of
$R \approx 9.6$~km and a central density $\varepsilon_c \approx
5\times10^{15}$~g~cm$^{-3}$ \cite{Oppy}. Beyond this point,
further increases in central density yield configurations of decreasing
mass and decreasing radius: these are unstable to radial perturbations
and would collapse to black holes if perturbed \cite{shapiro_teukolsky_1983}. Only
the branch of the $M$--$R$ curve with $dM/d\varepsilon_c > 0$ corresponds
to stable stellar configurations.

The termination of the stable branch at $M_{\mathrm{OV}}$ is a direct
consequence of the softening of the equation of state discussed in
Section~\ref{sec:tov}: as the central density
approaches and exceeds $\varepsilon_c^{\mathrm{max}}$, the neutrons
become increasingly ultra-relativistic and the equation of state
asymptotes toward $P \propto \varepsilon^{4/3}$, at which point the
three general relativistic correction terms in the TOV equation
collectively overwhelm the pressure support.


\subsection{Comparison with Observed Neutron Stars}

The free Fermi gas model can be assessed quantitatively by comparing its
$M$--$R$ predictions against well-measured neutron stars. Table~\ref{tab:neutron-stars}
lists a representative selection of neutron stars with precise mass
and/or radius measurements drawn from pulsar timing and X-ray
observations.

\begin{table}[h]
    \centering
    \caption{Selected neutron stars with well-measured masses and radii,
        compared to the Oppenheimer--Volkoff limit for the free
        neutron Fermi gas. Radii are from NICER X-ray timing.}

    \label{tab:neutron-stars}
    \begin{tabular}{lccc}
        \toprule
        Object & Mass ($M_\odot$) & Radius (km) & Reference \\
        \midrule
        PSR J0740+6620$^\dagger$  & $2.08 \pm 0.07$ & $12.39^{+1.30}_{-0.98}$ & \cite{Fonseca2021, Riley2021} \\
        PSR J0030+0451            & $1.34^{+0.15}_{-0.16}$ & $12.71^{+1.14}_{-1.19}$ & \cite{Riley2019} \\
        PSR J0437-4715            & $1.418 \pm 0.037$ & $11.36^{+0.95}_{-0.63}$ & \cite{Choudhury2024} \\
        PSR J1614-2230$^\dagger$  & $1.908 \pm 0.016$ & --- & \cite{Demorest2010} \\
        \midrule
        OV limit (free Fermi gas) & $\approx 0.70$ & $\approx 9.6$ & \cite{Oppy} \\
        \bottomrule
    \end{tabular}
\end{table}

Several features of Table~\ref{tab:neutron-stars} are immediately
striking. First, every observed neutron star has a mass that exceeds
$M_{\mathrm{OV}} \approx 0.7\,M_\odot$, with the most massive known
example (from this list), PSR~J0740+6620, reaching $2.08\,M_\odot$,nearly three times the OV limit. Second, the observed radii cluster around 11--13~km,
systematically larger than the $\approx 9.6$~km radius at the OV
maximum. Both discrepancies point in the same direction: the true
equation of state is stiffer than the free Fermi gas prediction,
meaning that at a given energy density the actual pressure exceeds what
a non-interacting neutron gas would produce. The physical origin of this is the repulsive component of the nuclear force at
supranuclear densities, which is absent from the free Fermi gas model
by construction.\cite{Lattimer2012}
